\documentclass{article}
\usepackage[T1]{fontenc}
\usepackage{lmodern}
\usepackage{graphicx}
\usepackage{amsmath,amssymb}
\usepackage[a4paper,margin=2cm]{geometry}
\usepackage[absolute,overlay]{textpos}
\setlength{\TPHorizModule}{1mm}
\setlength{\TPVertModule}{1mm}
\usepackage[indonesian]{babel}
\usepackage{hyperref}
\setlength{\emergencystretch}{2em}
% unit: Halaman 94

\begin{document}

% === Halaman 94 (python/native) ===

94

• Nilai terbesar yang dapat muncul untuk merepresentasikan blok: 25252525 

(ZZZZ),  maka 25252525 dapat digunakan sebagai modulus n. 

• Nilai m yang relatif prima dengan 25252525, misalnya 21035433,

• b dipilih antara 1 dan 25252525, misalnya 23210025. 

• Fungsi enkripsi menjadi:


C  21035433P + 23210025 (mod 25252525)

• Fungsi dekripsi, setelah dihitung, menjadi


P  5174971 (C – 23210025)  (mod 25252525)

\end{document}
